\begin{abstract}
% Agentic evolution has emerged as a powerful paradigm for improving programs, workflows, and scientific solutions.
% It iteratively generates candidate artifacts, evaluates them against a task objective, and uses the resulting feedback to guide subsequent evolution.
% Existing methods typically instantiate this paradigm either through fixed hand-designed procedures that are modular but rigid, or through general-purpose agents that flexibly integrate feedback but can drift as context grows, leaving long-horizon search vulnerable to local optima.
% %
% Both routes share a deeper limitation: the long-horizon evolution accumulates candidates, feedback, traces, and failures over time, but lacks a stable interface for organizing this evidence and revising the mechanism that drives future evolution.
% We therefore formulate \textbf{agentic evolution as an interactive environment}, where the accumulated evolution context becomes a process-level state.
% A meta-agent acts on this environment not by directly generating the next candidate, but by editing the mechanism that controls how future evolution proceeds.
% We introduce \method, a harnessed framework for meta-editing agentic evolution.
% It standardizes the evolution environment and provides a unified interface for observing accumulated evidence and editing the mechanism that drives future evolution.
% This lets \method revise both hand-designed procedures and agent operating contexts, reducing local-optimum risks in long-horizon evolution.
% Empirical evaluations on agentic and reasoning benchmarks show that \method outperforms 5 evolution baselines, achieving a \textbf{26\%} relative improvement over the strongest baseline.
% Furthermore, across 3 open-ended optimization tasks, \method outperforms 4 evolution baselines and achieves state-of-the-art performance under the same iteration budget.

% Agentic evolution has emerged as a powerful paradigm for improving programs, workflows, and scientific solutions by iteratively generating candidates, evaluating them, and using feedback to guide future search.
% However, existing methods are typically instantiated either as fixed hand-designed procedures that are modular but rigid, or as general-purpose agents that are flexible but prone to drift in long-horizon evolution.
% Both forms accumulate rich evidence over time, including candidates, feedback, traces, and failures, yet lack a stable interface for organizing this evidence and revising the search mechanism itself.
% We address this limitation by formulating \textbf{agentic evolution as an interactive environment}, where the accumulated evolution context serves as a process-level state.
% We introduce \method, a harnessed meta-editing framework in which a meta-agent observes this state and edits the procedure or agent context that controls future evolution, rather than directly proposing the next candidate.
% This unified interface enables \method to steer both procedure-based and agent-based evolution, improving evidence reuse and reducing the risk of long-horizon stagnation.
% Empirical evaluations on agentic and reasoning benchmarks show that \method outperforms five evolution baselines, achieving a \textbf{26\%} relative improvement over the strongest baseline.
% Across three open-ended optimization tasks, \method further outperforms four evolution baselines and achieves state-of-the-art performance under the same iteration budget.


Agentic evolution has emerged as a powerful paradigm for improving programs, workflows, and scientific solutions by iteratively generating candidates, evaluating them, and using feedback to guide future search.
However, existing methods are typically instantiated either as fixed hand-designed procedures that are modular but rigid, or as general-purpose agents that flexibly integrate feedback but can drift in long-horizon evolution.
Both forms accumulate rich evidence over time, including candidates, feedback, traces, and failures, yet lack a stable interface for organizing this evidence and revising the mechanism that drives future evolution.
We address this limitation by formulating \textbf{agentic evolution as an interactive environment}, where the accumulated evolution context serves as a process-level state.
We introduce \method, a harnessed meta-editing framework in which a meta-agent observes this state and acts not by directly proposing the next candidate, but by editing the procedure or agent context that controls future evolution.
This unified interface enables \method to steer both procedure-based and agent-based evolution, making accumulated evidence actionable for long-horizon search.
Empirical evaluations on agentic and reasoning benchmarks show that \method outperforms five evolution baselines, achieving a \textbf{26\%} relative improvement over the strongest baseline.
Across three open-ended optimization tasks, \method further outperforms four evolution baselines and achieves state-of-the-art performance under the same iteration budget.
Code is available at \href{https://github.com/FoundationAgents/AEvo}{FoundationAgents/AEvo}.
\end{abstract}
